\section{Response-Aware Marginal-Utility Routing}

R-MUR estimates the marginal utility of a candidate with a fixed predictor. It
screens candidates with cached representations, evaluates the predictor only on
a shortlist, and estimates the marginal from the cached context together with
the predictor's response.

\subsection{Marginal-utility estimation}

Every candidate is represented by a cached feature vector $h_j$ computed from
its context, and the query and the current selected set are represented by
$h_Q$ and $h_A$. The selected-set encoder is permutation invariant, so the
representation depends on the set and not on the order in which its elements
were added. A cached model maps these representations to a utility score
\begin{equation}
 \widehat m_C(j\mid A)=g_C(h_Q,h_A,h_j),
 \label{eq:cached}
\end{equation}
which is inexpensive because it never evaluates the predictor. It is also
incomplete: cached features describe the content of a candidate, not how the
predictor reacts to it.

\subsection{Response-aware routing}

The cached score forms a shortlist
\begin{equation}
 \mathcal Q_A=\operatorname*{TopQ}_{j\notin A}\widehat m_C(j\mid A),
 \qquad |\mathcal Q_A|=q\ll K,
 \label{eq:shortlist}
\end{equation}
which bounds the number of predictor evaluations per decision to $q$ out of
$K$ candidates. For each retained candidate the predictor is evaluated with and
without that candidate, and a compact response summary is formed,
\begin{equation}
 r_{A,j}=\psi\!\left(f_A(x),f_{A\cup\{j\}}(x)\right).
 \label{eq:response}
\end{equation}
The summary contains the level of the current prediction, the level after
adding the candidate, and the mean, standard deviation, mean absolute value,
normalized Euclidean norm, and maximum absolute value of the change. These
statistics describe how the candidate changes the prediction, which is the
quantity that enters the marginal utility through
Eq.~\eqref{eq:squared-marginal}. The cached representations are the anchor history, the candidate history, and
the selected candidate histories with a validity mask. The response-aware model
combines them with the cached representations,
\begin{equation}
 \widehat m_R(j\mid A)=g_R(h_Q,h_A,h_j,r_{A,j}),
 \qquad j\in\mathcal Q_A .
 \label{eq:rmur}
\end{equation}

\subsection{Training and sequential inference}

The model is trained on states sampled from the training episodes. A state is
a uniformly random subset of the candidate pool whose size is drawn uniformly
from $0$ to $B-1$, so training covers empty, partial and nearly complete
selected sets; deployment states follow the router's own greedy trajectory
instead, and this difference between the training and deployment state
distributions is the main approximation in the design. At each state the
counterfactual loss difference of Eq.~\eqref{eq:target} provides the observed
marginal for every candidate, so the labels are state-conditioned rather than
standalone, and one state supplies $K-|A|$ regression targets together with all
candidate pairs for the ranking term.

The encoder is a permutation-invariant set encoder of depth two with hidden and
set width 64, and the utility head is a two-layer network of width 64. Training
uses AdamW with learning rate $10^{-3}$ and weight decay $10^{-4}$ for 20
epochs with batches of 128 states; each target-run contributes 16{,}000
training states, four per episode, with 16 counterfactual labels each. The objective combines regression with a pairwise ranking term,
\begin{equation}
 \mathcal L
 =
 \mathcal L_{\mathrm{Huber}}
 +
 \beta
 \sum_{i,j:\,\widetilde m_i>\widetilde m_j}
 |\widetilde m_i-\widetilde m_j|
 \log\!\left(
 1+\exp\!\left[-(\widehat m_i-\widehat m_j)\right]
 \right),
 \label{eq:loss}
\end{equation}
where the sum is taken over pairs within the same selection state and the
weight is the utility gap of the pair. The gap weight assigns greater
importance to ranking errors with larger utility consequences, while giving
little weight to nearly tied candidates, and the regression term keeps the
utility scale calibrated on states without a dominant alternative. The
response-aware scorer is trained over candidate states sampled from the full
training distribution, where every candidate of a state carries a label and
enters the pairwise term. At inference the scorer is evaluated only on the
shortlist produced by the cached scorer.

At deployment the router starts from the empty set and repeats the following
steps: score all remaining candidates with the cached model, form the
shortlist, evaluate the predictor on the shortlist, rerank it with the
response-aware model, and add the highest-scoring candidate when its estimated
marginal utility is positive. The procedure stops when no shortlisted candidate
has a positive estimated utility or when the context budget is reached, and it
is greedy: each decision maximizes the estimated one-step marginal utility
within the shortlist. Section~\ref{sec:analysis} analyzes the regret that
screening and reranking introduce into this two-stage decision.
Algorithm~\ref{alg:rmur} states the procedure.

\begin{algorithm}[t]
\caption{Response-Aware Marginal-Utility Routing}
\label{alg:rmur}
\begin{algorithmic}[1]
\Require anchor $S_0$, query $x$, candidate pool $\candidates$, budget $B$, shortlist size $q$, fixed predictor $f$
\State $A \gets \emptyset$
\While{$|A| < B$}
  \State $\widehat m_C(j\mid A) \gets g_C(h_Q,h_A,h_j)$ for all $j\notin A$ \Comment{cached screening}
  \State $\mathcal Q_A \gets \operatorname*{TopQ}_{j\notin A}\widehat m_C(j\mid A)$
  \ForAll{$j \in \mathcal Q_A$}
    \State $r_{A,j} \gets \psi(f_A(x), f_{A\cup\{j\}}(x))$ \Comment{predictor response}
    \State $\widehat m_R(j\mid A) \gets g_R(h_Q,h_A,h_j,r_{A,j})$
  \EndFor
  \State $j^\star \gets \argmax_{j\in\mathcal Q_A}\widehat m_R(j\mid A)$
  \If{$\widehat m_R(j^\star\mid A) \le 0$}
    \State \textbf{break} \Comment{stop: no shortlisted candidate is estimated useful}
  \EndIf
  \State $A \gets A\cup\{j^\star\}$
\EndWhile
\State \Return $A$
\end{algorithmic}
\end{algorithm}
